%!TEX TS-program = xelatex
%!TEX TS-options = -shell-escape
%!TEX root = ../Topo2_WS15/topologie_2.tex
\RequirePackage{fix-cm} 
\documentclass[a4paper, twoside, headsepline, index=totoc,toc=listof, fontsize=10pt, cleardoublepage=empty, headinclude, DIV=12, BCOR=5mm, titlepage,draft]{scrartcl}
\usepackage{scrtime} % KOMA, Uhrzeit ermoeglicht

\usepackage{etoolbox}
\usepackage{letltxmacro}
\usepackage{ifthen}

%--Farbdefinitionen
%-- muss vor tikz geladen werden
\usepackage[usenames, table, x11names]{xcolor}
\definecolor{dark_gray}{gray}{0.45}
\definecolor{light_gray}{gray}{0.6}
\usepackage[final]{graphicx}

%--Zum Zeichnen
%-- muss vor polyglossia bzw. babel geladen werden
\usepackage{tikz}
\usepackage{tikz-cd}
\usetikzlibrary{external}
\tikzset{>=latex}
\usetikzlibrary{shapes,arrows.meta,intersections}
\usetikzlibrary{calc,3d}
\usetikzlibrary{decorations.pathreplacing,decorations.markings, decorations.pathmorphing}
\usetikzlibrary{angles}

%-- Konfiguration von tikzexternalize
\tikzexternalize[prefix=tikz/,up to date check=diff]
\pgfkeys{/pgf/images/include external/.code=\includegraphics{#1}}
\tikzset{external/system call={xelatex \tikzexternalcheckshellescape -halt-on-error -interaction=batchmode --shell-escape -jobname "\image" "\texsource"}}


%-- tikzexternalize fuer tikzcd deaktivieren, da inkompatibel
\AtBeginEnvironment{tikzcd}{\tikzexternaldisable}
\AtEndEnvironment{tikzcd}{\tikzexternalenable}

%-- um Inkompatibilitaeten von quotes und polyglossia bzw. babel zu vermeiden
\tikzset{
  every picture/.append style={
    execute at begin picture={\shorthandoff{"}},
    execute at end picture={\shorthandon{"}}
  }
}
\usetikzlibrary{quotes}
\usepackage{pgfplots}
\usepgfplotslibrary{colormaps}
% \pgfplotsset{compat=1.12}



%-- Mathesymbole etc.
\usepackage{mathtools} % beinhaltet amsmath
\mathtoolsset{showonlyrefs, centercolon,showmanualtags}
\newtagform{brackets}[\textbf]{[}{]}
\usetagform{brackets}
\usepackage{fix-cm}
% \usepackage{wasysym} % zusätzliche Symbole
\usepackage{amssymb} % zusätzliche Symbole
% \usepackage{stmaryrd} % für Blitz
\usepackage{nicefrac} % schräge Brüche
\usepackage{faktor}
\newcommand{\Faktor}[1]{\faktor[\textstyle]{#1}}
\usepackage{xfrac}
\usepackage{cancel}
\usepackage{mathdots} % Verbesserung von Punkten wie zB \ldots
\usepackage[bb=px]{mathalfa}

%-- einzelne Symbole aus dem mathx-package
% \DeclareFontFamily{U}{mathx}{\hyphenchar\font45}
% \DeclareFontShape{U}{mathx}{m}{n}{<-> mathx10}{}
% \DeclareSymbolFont{mathx}{U}{mathx}{m}{n}
% \DeclareMathSymbol{\bigplus}{1}{mathx}{"90}

%-- 
\usepackage{silence}
\WarningFilter{latexfont}{Size substitutions with differences}
\WarningFilter{latexfont}{Font shape `U/bbold/m/n' in size}

\DeclareSymbolFont{bbold}{U}{bbold}{m}{n}
\DeclareSymbolFontAlphabet{\mathbbold}{bbold}
\newcommand{\ind}{\mathbbold{1}} % charakteristische-Funktion-Eins

\def\mathul#1#2{\color{#1}\underline{{\color{black}#2}}\color{black}} %farbiges Untersteichen im Mathe-Modus
\renewcommand{\le}{\leqslant}
\renewcommand{\ge}{\geqslant}

%-- Underbrace als Befehl in LaTeX-Syntax (und ohne Spacingprobleme mit nachfolgenden Operatoren...)
\newcommand{\Underbrace}[2]{{\underbrace{#1}_{#2}}}
\newcommand{\Underbracket}[2]{{\underbracket[0.7pt][2pt]{#1}_{#2}}}

%-- Alles was mit Schrift und XeTeX zu tun hat
\usepackage{palatino}
\usepackage[euler-digits]{eulervm}
\usepackage[no-math]{fontspec}
\usepackage{polyglossia} % moderner babel-ersatz
\setmainlanguage[spelling=new,babelshorthands=true]{german}
\shorthandoff{"}
\setotherlanguage{english}
\defaultfontfeatures{Mapping=tex-text, WordSpace={1.2}, Ligatures={Required,Common,Contextual}} %

%-- Alle hier aufgeführten Schriftarten sind in jeder TeXlive-Installation vorhanden!!
%-- Hauptschriftart festlegen
% \setmainfont{texgyrepagella}[Extension=.otf,UprightFont=*-regular,BoldFont=*-bold,ItalicFont=*-italic,BoldItalicFont=*-bolditalic,ItalicFeatures={Style=Historic}]
\setmainfont{TeXGyrePagellaX}[Extension=.otf,UprightFont=*-Regular,BoldFont=*-Bold,ItalicFont=*-Italic,BoldItalicFont=*-BoldItalic,ItalicFeatures={Style=Historic},Ligatures={Required,Common,Contextual,Historic}]
% \setmainfont[BoldFont={* Bold}, ItalicFont={* Light Italic},SmallCapsFont={Linux Libertine O}, SmallCapsFeatures={Letters=SmallCaps}]{Source Sans Pro}


%-- Serifenlose Schriftart festlegen
% \setsansfont{Roboto}[Scale=MatchUppercase,Extension=.ttf,UprightFont=*-Regular,BoldFont=*-Bold,ItalicFont=*-RegularItalic,BoldItalicFont=*-MediumItalic]
% \setsansfont{RobotoSlab}[Scale=MatchUppercase,Extension=.ttf,UprightFont=*-Regular,BoldFont=*-Bold,ItalicFont=*-Thin,SmallCapsFont=*-Regular]
% \setsansfont{FiraSans}[Scale=MatchUppercase,Extension=.otf, UprightFont=*-Regular, BoldFont=*-SemiBold, ItalicFont=*-RegularItalic, BoldItalicFont=*-MediumItalic]
% \setsansfont{SourceSansPro}[Scale=MatchUppercase,Extension=.otf,UprightFont=*-Regular,BoldFont=*-Semibold,ItalicFont=*-LightIt,BoldItalicFont=*-SemiboldIt]
% \setsansfont{texgyreadventor}[Scale=MatchUppercase, Extension=.otf, UprightFont=*-regular, BoldFont=*-bold, ItalicFont=*-italic,BoldItalicFont=*-bolditalic]
% \setsansfont{Kurier}[Scale=MatchUppercase, Extension=.otf, UprightFont=*Medium-Regular, BoldFont=*Heavy-Regular, ItalicFont=*Medium-Italic,BoldItalicFont=*Heavy-Italic]
% \setsansfont{Iwona}[Scale=MatchUppercase, Extension=.otf, UprightFont=*Medium-Regular, BoldFont=*Heavy-Regular, ItalicFont=*Medium-Italic,BoldItalicFont=*Heavy-Italic]
% \setsansfont{LinBiolinum}[Scale=MatchUppercase, Extension=.otf, UprightFont=*_R, BoldFont=*_RB, ItalicFont=*_RI,BoldItalicFont=*_RBO]
% \setsansfont{texgyrebonum}[Scale=MatchUppercase,Extension=.otf, UprightFont=*-regular, BoldFont=*-bold, ItalicFont=*-italic, BoldItalicFont=*-bolditalic]
\setsansfont{texgyreadventor}[Scale=MatchUppercase,Extension=.otf, UprightFont=*-regular, BoldFont=*-bold, ItalicFont=*-italic, BoldItalicFont=*-bolditalic]



% \newfontfamily{\geometric}{texgyreheros}[Scale=MatchLowercase, Extension=.otf, UprightFont=*-regular, BoldFont=*-bold, ItalicFont=*-italic]
% \addtokomafont{disposition}{\geometric}
\setmonofont{SourceCodePro}[Scale=0.9,Extension=.otf,UprightFont=*-Regular, BoldFont=*-Semibold, ItalicFont=*-Light]
\usepackage{xltxtra}
\usepackage{fontawesome}
\usepackage[final]{microtype}


\usepackage[draft=false]{scrlayer-scrpage} % wie fancyhdr, nur optimiert auf KOMA-Skript, leicht andere Syntax

\flushbottom


%--Mixed
\usepackage[shortlabels,inline]{enumitem}
\setlist[itemize,1]{label=\faCaretRight}
\setlist[enumerate]{font=\bfseries}
\setlist[description]{font=\normalfont\bfseries}
\usepackage[german=quotes]{csquotes}
\usepackage{makeidx}
\usepackage{wrapfig}
\usepackage{float}
\usepackage[margin=10pt, font=small, labelfont={sf, bf}, format=plain, indention=1em]{caption}
\captionsetup[wrapfigure]{name=Abb. }
\usepackage{stackrel}


%--Unterstreichung
\usepackage[normalem]{ulem}
\setlength{\ULdepth}{1.8pt}

%--Indexverarbeitung
\newcommand{\bet}[1]{\emph{#1}}
\newcommand{\Index}[1]{\bet{#1}\index{#1}}
\makeindex
\setindexpreamble{{\noindent\sffamily\small Die \emph{Seitenzahlen} sind mit Hyperlinks versehen und somit anklickbar}  \par \bigskip}
\renewcommand{\indexpagestyle}{scrheadings}


%--Marginnote & todonotes
\deffootnote[1.5em]{1.5em}{1.5em}{\textsuperscript{\thefootnotemark}\ }
\usepackage[fulladjust]{marginnote}
\renewcommand*{\marginfont}{\color{dark_gray} \itshape \footnotesize}
\usepackage[textsize=small]{todonotes}
\usepackage{ragged2e}
\renewcommand*{\raggedleftmarginnote}{\RaggedLeft}
\renewcommand*{\raggedrightmarginnote}{\RaggedRight}

%--Todonotes mit tikz/externalize kompatibel machen
\LetLtxMacro{\oldtodo}{\todo}
\renewcommand{\todo}[2][]{\tikzexternaldisable\oldtodo[#1]{#2}\tikzexternalenable}
\LetLtxMacro{\oldmissingfigure}{\missingfigure}
\renewcommand{\missingfigure}[2][]{\tikzexternaldisable\oldmissingfigure[{#1}]{#2}\tikzexternalenable}

%--Konfiguration von Hyperref pdfstartview=FitH, 
\usepackage[hidelinks, pdfpagelabels,  bookmarksopen=true, bookmarksnumbered=true, linkcolor=black, urlcolor=SkyBlue2, plainpages=false,pagebackref, citecolor=black, hypertexnames=true, pdfauthor={Jannes Bantje}, pdfborderstyle={/S/U}, linkbordercolor=SkyBlue2, colorlinks=false,final]{hyperref}


\newcommand{\appendLink}[1]{#1\,\faExternalLink}
\newcommand{\hrefsym}[2]{\href{#1}{\texttt{\appendLink{#2}}}}
\newcommand{\hrefsymmail}[2]{\href{#1}{\texttt{\faEnvelopeO\,#2}}}
\renewcommand{\url}[1]{\hrefsym{#1}{\nolinkurl{#1}}}



% -- QR-Codes
\usepackage{qrcode} %-- hinter hyperref laden!

%--Römische Zahlen
\newcommand{\RM}[1]{\MakeUppercase{\romannumeral #1{}}}

%%--Abkürzungen etc.
\newcommand{\light}{\text{\Large $\lightning$}}

%-- Definitionen von weiteren Mathe-Befehlen

\DeclareMathOperator{\im}{im} %Bild
\DeclareMathOperator{\id}{id} %identische Abbildung
\DeclareMathOperator{\conj}{conj} %Konjugation
\DeclareMathOperator{\sgn}{sgn} %Signum
\DeclareMathOperator{\End}{End} %Endomorphismen
\DeclareMathOperator{\Hom}{Hom} %Homomorphismen
\DeclareMathOperator{\Iso}{Iso} %Isomorphismen
\DeclareMathOperator{\Span}{span} %Span

\DeclareMathOperator{\grad}{grad} %Gradient
\DeclareMathOperator{\dive}{div} %Gradient
\DeclareMathOperator{\rot}{rot} %Rotation


\DeclareMathOperator{\Char}{char} %Charakteristik
\DeclareMathOperator{\Aut}{Aut} %Automorphismen

% -- Zum Finetuning von Befehlen
\makeatletter
\newcommand{\raisemath}[1]{\mathpalette{\raisem@th{#1}}}
\newcommand{\raisem@th}[3]{\raisebox{#1}{$#2#3$}}
\makeatother

\newcommand{\sing}{{\raisemath{1.1pt}{\scriptscriptstyle\mathrm{sing}}}}
\newcommand{\pt}{\mathrm{pt}}
\newcommand{\op}{\mathrm{op}}
\DeclareMathOperator{\Sp}{Sp}
\DeclareMathOperator{\Keg}{Keg}
\newcommand{\slashedi}{i\hspace{-3.5pt}/}
\newcommand{\cupp}{\smallsmile}
\newcommand{\capp}{\smallfrown}
\DeclareMathOperator*{\colim}{colim}


\newcommand{\D}{\ensuremath{\mathrm{D}\mkern-1.5mu}}
\newcommand{\mathd}{\ensuremath{\mathrm{d}\mkern-1.0mu}}
\newcommand{\Tmap}{\ensuremath{\mathrm{T}\mkern-0.85mu}}
\newcommand{\opL}{\ensuremath{\mathrm{L}\mkern-0.6mu}}


\DeclareMathOperator{\supp}{supp} %Träger


% -- Befehle für FunkAna und Operatoralgebren
\newcommand{\tr}{\mathrm{tr}}
\newcommand{\w}{\mathrm{w}}
\newcommand{\sa}{\mathrm{sa}}
\newcommand{\so}{\mathrm{s.o.}}
\newcommand{\weakT}[1]{\ensuremath{\mathcal{T}_{#1}^{\mkern+1.0mu\text{\raisebox{0.4ex}{$\mathrm{w}$}}}}}
\newcommand{\weakTstar}[1]{\ensuremath{\mathcal{T}_{#1}^{\mkern+1.0mu\text{\raisebox{0.4ex}{$\mathrm{w}$}}^*}}}
\newcommand{\finSub}{\subset\mkern-0.7mu \subset}
\DeclareMathOperator{\Inv}{Inv}
\DeclareMathOperator{\dist}{dist}
\newcommand{\simm}{{\hspace{-1.6pt}\raisemath{0.5pt}{\sim}}}
\DeclareMathOperator{\ev}{ev}
\DeclareMathOperator{\Alg}{Alg}


% -- Kategorien
\newcommand{\TOP}{\textsc{Top}}
\newcommand{\HTOP}{\textsc{HTop}}
\newcommand{\VR}{\textsc{VR}}
\newcommand{\MOD}{\textsc{Mod}}
\newcommand{\MONOIDE}{\textsc{Monoide}}
\newcommand{\MENGEN}{\textsc{Mengen}}
\newcommand{\MAN}{\textsc{Man}}
\newcommand{\GRUPPEN}{\textsc{Gruppen}}
\newcommand{\ABELGRUPPEN}{\textsc{Abel.Gruppen}}
\newcommand{\KAT}{\textsc{Kat}}
\newcommand{\FUN}{\textsc{Fun}}
\newcommand{\SIMP}{\textsc{Simp}}
\newcommand{\VEKT}{\textsc{Vekt}}

\DeclareMathOperator{\cell}{cell}
\DeclarePairedDelimiter{\homologieklasse}{\llbracket}{\rrbracket}
\newcommand{\rand}[1]{\ensuremath{\partial^{\scriptscriptstyle #1}}}


% -- theorem packages
\usepackage{amsthm}
\usepackage{thmtools,thm-restate}
\renewcommand{\listtheoremname}{Übersicht aller Aussagen}

% -- Theoreme als PDF-Lesezeichen
\usepackage{bookmark}
\bookmarksetup{open,numbered}
\makeatletter
\newcommand*{\theorembookmark}{%
  \bookmark[
    dest=\@currentHref,
    rellevel=1,
    keeplevel,
  ]{%
    \thmt@thmname\space\csname the\thmt@envname\endcsname
    \ifx\thmt@shortoptarg\@empty
    \else
      \space(\thmt@shortoptarg)%
    \fi
  }%
}   
\makeatother

% -- Definition der einzelnen Umgebungen
\declaretheoremstyle[%
	headfont=\sffamily\bfseries,
	notefont=\normalfont\sffamily\scshape,
	bodyfont=\normalfont,
	headformat=\NUMBER\ \NAME\NOTE,
	headpunct=.,
	postheadspace=1em,
	spaceabove=15pt,spacebelow=10pt,
	postheadhook=\theorembookmark]%
{mainstyle}
\declaretheoremstyle[%
	headfont=\sffamily\bfseries,
	notefont=\normalfont\sffamily\scshape,
	bodyfont=\normalfont,
	headformat=\NUMBER\NAME\NOTE,
	headpunct=.,
	postheadspace=1em,
	spaceabove=15pt,spacebelow=10pt,
	postheadhook=\theorembookmark]%
{mainstyle_unnumbered}
\declaretheoremstyle[%
	headfont=\sffamily\bfseries,
	notefont=\normalfont\sffamily\scshape,
	bodyfont=\normalfont,
	headformat=swapnumber,
	headpunct=.,
	postheadspace=1em,
	spaceabove=15pt,spacebelow=10pt,
	postheadhook=\theorembookmark,
	qed=\qedsymbol]%
{mainstyleB}
\declaretheorem[name=Definition,parent=section,style=mainstyle]{definition}
\declaretheorem[name=Definition,numbered=no,style=mainstyle_unnumbered]{definition*}
\declaretheorem[name=Theorem,sharenumber=definition,style=mainstyle]{theorem}
\declaretheorem[name=Theorem,numbered=no,style=mainstyle_unnumbered]{theorem*}
\declaretheorem[name=Proposition,sharenumber=definition,style=mainstyle]{proposition}
\declaretheorem[name=Lemma,sharenumber=definition,style=mainstyle]{lemma}
\declaretheorem[name=Satz,sharenumber=definition,style=mainstyle]{satz}
\declaretheorem[name=Korollar,sharenumber=definition,style=mainstyle]{korollar}
\declaretheorem[name=Korollar,sharenumber=definition,style=mainstyleB]{korollarB}
\declaretheorem[name=Frage,numbered=no,style=mainstyle_unnumbered]{frage}
\declaretheorem[name=Erinnerung,sharenumber=definition,style=mainstyle]{erinnerungA}


% -- Beweise
\declaretheoremstyle[headfont=\bfseries\scshape,bodyfont=\normalfont,headpunct=:,postheadspace=1em,spacebelow=12pt,spaceabove=2pt,qed=\qedsymbol]{beweise}
\declaretheoremstyle[headfont=\sffamily\bfseries,bodyfont=\normalfont,headpunct=:,postheadspace=1em,spacebelow=10pt,spaceabove=10pt]{bemerkungen}
\declaretheorem[name=Beweis,numbered=no,style=beweise]{beweis}
\declaretheorem[name=Bemerkung,sharenumber=definition,style=mainstyle]{bemerkung}
\declaretheorem[name=Beispiel,sharenumber=definition,style=mainstyle]{beispiel}
\declaretheorem[name=Übung,numbered=no,style=bemerkungen]{uebung}
\declaretheorem[name=Erinnerung,numbered=no,style=bemerkungen]{erinnerung}


%--Skalarprodukt
\DeclarePairedDelimiterX\sprod[2]{\langle}{\rangle}{#1\,\delimsize\vert\,#2}
\DeclarePairedDelimiterX\skal[2]{\langle}{\rangle}{#1\,,\,#2}

%--Betrag, Gaußklammer
\DeclarePairedDelimiter{\abs}{\lvert}{\rvert}
\DeclarePairedDelimiter{\floor}{\lfloor}{\rfloor}
\DeclarePairedDelimiter{\ceil}{\lceil}{\rceil}

%--Norm
\DeclarePairedDelimiter\norm{\Vert}{\Vert}


%--Umklammern
\DeclarePairedDelimiter\enbrace{(}{)}
\DeclarePairedDelimiter\benbrace{[}{]}
\DeclarePairedDelimiter\lenbrace{<}{>}
\newcommand{\ssbrace}[1]{{\scriptscriptstyle\enbrace{#1}}}

%--Mengen
\newcommand\SetSymbol[1][]{\nonscript\:#1\vert\allowbreak\nonscript\:\mathopen{}}
\providecommand\given{} % to make it exist
\DeclarePairedDelimiterX\set[1]\{\}{\renewcommand\given{\SetSymbol[\delimsize]}#1}

%--Abbildungsdefinition
\newcommand{\mapdef}[5]{%
	\[
		\begin{array}{rcl}
			\textstyle #1 &\xrightarrow{\minwidthbox{#5}{2em}} & \textstyle #2 \\[0.5ex]
			\textstyle #3 &\xmapsto{\minwidthbox{\mbox{ }}{2em}} & \textstyle #4
		\end{array}
	\]
}

%--modifiziertes Stackrel
\newcommand{\StackText}[2]{\stackrel{\mbox{\scriptsize #1}}{#2}}

%--Befehl für eine Box mit Mindestweite
\DeclareRobustCommand{\minwidthbox}[2]{%
  \ifmmode
    \expandafter\mathmakebox
  \else
    \expandafter\makebox
  \fi
  [\ifdim#2<\width\width\else#2\fi]{#1}%
}

%--Befehl zum Einkreisen
\newcommand*\circled[1]{\tikzexternaldisable\tikz[baseline=(char.base)]{\node[shape=circle,draw,inner sep=2pt] (char) {#1};}\tikzexternalenable}


%--Differential
\newcommand{\diff}[2]{\ensuremath{\frac{{\partial #1}}{{\partial #2}} }}
\newcommand{\diffd}[2]{\ensuremath{\frac{\mathd #1}{\mathd #2} }}

%--Inhaltsverzeichnis
\usepackage[tocindentauto]{tocstyle}
\usetocstyle{KOMAlike}	
\shorthandon{"}
		
%--Punkte (als letztes laden)
\usepackage{ellipsis}

\newcommand{\fach}{Topologie \RM{1}.}
\newcommand{\semester}{WiSe 2014/2015}
\newcommand{\homepage}{https://wwwmath.uni-muenster.de/reine/u/topos/lehre/WS2014-2015/Topologie1/Topologie1.html}
\publishers{\scalebox{11}{\Huge$H_*$}}
\newcommand{\prof}{Prof.\ Dr.\ Arthur Bartels}
\providecommand{\verfasser}{Jannes Bantje}

%--Konfiguration von scrheadings
\setheadsepline{1pt}[\color{light_gray}]
\pagestyle{scrheadings}
%\fancyhf[H, F]{}
\clearscrheadfoot

\providecommand{\shortFach}{\fach}
\lehead{ \includegraphics[height=0.5 cm,keepaspectratio]{../!config/Bilder/Logo_WWU_Muenster_light_gray.pdf}}
\rehead{\rule{0cm}{0.5cm}\footnotesize \sffamily \color{light_gray} \verfasser -- Mitschrift \shortFach}
\lohead{\rule{0cm}{0.5cm} \footnotesize \sffamily \color{light_gray} Stand: \today \; \thistime[:]}
\rohead{\includegraphics[height=0.5 cm,keepaspectratio]{../!config/Bilder/fb10logo_gray.pdf}}


\ofoot[{ \color{dark_gray} \LARGE \sffamily \thepage}]{{ \color{dark_gray} \LARGE \sffamily \thepage}} %hier wir auch der plain Stil bearbeitet!
\automark{section}
\ifoot{ \color{dark_gray} \small \leftmark}

%--Metadaten
\author{\verfasser}
\titlehead{\includegraphics[height=1.5cm, keepaspectratio]{../!config/Bilder/Logo_WWU_Muenster.pdf}%
\hfill \includegraphics[height=1.3cm, keepaspectratio]{../!config/Bilder/fb10logo.pdf}}
\title{Skript \fach}
\subtitle{Mitschrift der Vorlesung  \enquote{\fach} von \prof}

\publishers{
Aktuelle Version verfügbar bei: \bigskip\\
\normalsize
\begin{minipage}[t]{0.45\textwidth}
	\centering{
	\qrcode[height=3.5cm, version=4]{https://github.com/JaMeZ-B/latex-wwu}
	\medskip\\
	\includegraphics[height=0.4cm, keepaspectratio]{../!config/Bilder/github_octo.pdf} 
	\includegraphics[height=0.4cm, keepaspectratio]{../!config/Bilder/GitHub_Logo.pdf} (inklusive Sourcecode)\\
	{\small\url{https://github.com/JaMeZ-B/latex-wwu}}}
\end{minipage}
\hfill
\begin{minipage}[t]{0.45\textwidth}
	\centering{
	\qrcode[height=3.5cm, version=4]{btsync://B6WH2DISQ5QVYIRYIEZSF4ZR2IDVKPN3I?n=latex_share}
	\medskip\\
	\includegraphics[height=0.4cm, keepaspectratio]{../!config/Bilder/bt_sync_logo.pdf} 
	{\large \textbf{Bittorrent} Sync} \\
	\texttt{B6WH2DISQ5QVYIRYIEZSF4ZR2IDVKPN3I}
	}
\end{minipage}
}


\begin{document}
\pagenumbering{Roman}
\maketitle
\begin{abstract}
\section*{Vorwort --- Mitarbeit am Skript}
Dieses Dokument ist eine Mitschrift aus der Vorlesung \enquote{\fach, \semester}, gelesen von \prof. Der Inhalt entspricht weitestgehend dem Tafelanschrieb. Für die
Korrektheit des Inhalts übernehme ich keinerlei Garantie! Für Bemerkungen und Korrekturen -- und seien es nur Rechtschreibfehler -- bin ich sehr dankbar. 
Korrekturen lassen sich prinzipiell auf drei Wegen einreichen: 
\begin{itemize}
	\item Persönliches Ansprechen in der Uni, Mails an \hrefsymmail{mailto:\mail}{\mail} (gerne auch mit annotieren PDFs) 
	\item \emph{Direktes} Mitarbeiten am Skript: Den Quellcode poste ich auf GitHub (siehe oben), also stehen vielfältige Möglichkeiten der Zusammenarbeit zur Verfügung:
	Zum Beispiel durch Kommentare am Code über die Website und die Kombination Fork + Pull Request. Wer sich verdient macht oder ein Skript zu einer Vorlesung, die 
	ich nicht besuche, beisteuern will, dem gewähre ich gerne auch Schreibzugriff.
	
	Beachten sollte man dabei, dass dazu ein Account bei \url{github.com} notwendig ist, der allerdings ohne Angabe von persönlichen Daten angelegt werden kann. 
	Wer bei GitHub (bzw. dem zugrunde liegenden Open-Source-Programm \enquote{\texttt{git}}) -- verständlicherweise -- Hilfe beim Einstieg braucht, dem helfe ich gerne 
	weiter. Es gibt aber auch zahlreiche empfehlenswerte Tutorials im Internet.\footnote{zB. \url{https://try.github.io/levels/1/challenges/1}, ist auf Englisch, aber dafür 
	interaktives LearningByDoing}
	\item \emph{Indirektes} Mitarbeiten: \TeX-Dateien per Mail verschicken. 
	
	Dies ist nur dann sinnvoll, wenn man einen ganzen Abschnitt ändern möchte (zB. einen alternativen Beweis geben), da ich die Änderungen dann per Hand einbauen muss! Ich freue mich aber auch über solche Beiträge!
\end{itemize}
\section*{Vorlesungshomepage}
{\centering 
\begin{minipage}[c][][c]{\textwidth}
	\centering \qrcode[height=3cm]{\homepage} \medskip\\
	\footnotesize \url{\homepage}
\end{minipage}
\par}

\section*{Inhalt der Vorlesung}
In der Vorlesung werden wir uns mit algebraischer Topologie beschäftigen. 
In der algebraischen Topologie werden topologische Fragen, wie zum Beispiel wann $\mathbb{R}^n$ und $\mathbb{R}^m$ homöomorph sind, in algebraische Fragen übersetzt. 
Ein wichtiges Hilfsmittel für solche Übersetzungen sind Homologiegruppen. 
Wir werden in der Vorlesung diese Homologiegruppen definieren und Hilfsmittel entwickeln, um Homologiegruppen zu berechnen. 
Es ergeben sich dann schöne topologische Anwendungen, wie zum Beispiel die topologische Invarianz der Dimension. 
\section*{Literatur}
\begin{itemize}
	\item \textsc{Bredon, G.E.}: \emph{Topology and Geometry}, Graduate Texts in Mathematics, Springer Verlag 1993
	\item \textsc{Hatcher, A.}: \href{http://www.math.cornell.edu/~hatcher/AT/ATpage.html}{\appendLink{\emph{Algebraic Topology}}}, Cambridge University Press, 2002
	\item \textsc{Lück, W.}: \emph{Algebraische Topologie - Homologie und Mannigfaltigkeiten}, Vieweg, 2005
	\item \textsc{tom Dieck, T.}: \emph{Topologie}, 2. Auflage, de Gruyter 2000
	\item \textsc{Waldhausen, F.}: \href{https://www.math.uni-bielefeld.de/~fw/ein.pdf}{\appendLink{\emph{Skript zur Einführung in die Topologie}}}, 
	\href{https://www.math.uni-bielefeld.de/~fw/at.pdf}{\appendLink{\emph{Skript zur Topologie}}}
\end{itemize}
\end{abstract}

\tableofcontents
\cleardoubleoddemptypage

\pagenumbering{arabic}
\setcounter{page}{1}
\setcounter{footnote}{0}

\section{Kategorien, Funktoren und natürliche Transformationen} % (fold)
\label{sec:1}

\begin{definition}[{name=[Kategorie]}]
	Eine \Index{Kategorie} $\mathcal{C}$ besteht aus:
	\begin{enumerate}[i),itemsep=0pt]
		\item Einer Klasse $\Obj(\mathcal{C})$. Die Elemente von $\Obj(\mathcal{C})$ heißen die \Index{Objekte} von $\mathcal{C}$.
		\item Zu je zwei Objekten $A,B$ aus $\mathcal{C}$ einer Menge $\Mor_{\mathcal{C}}(A,B)$. Die Elemente von $\Mor_\mathcal{C}(A,B)$ heißen \Index{Morphismen} von $\mathcal{C}$.
		\item Zu je drei Objekten $A,B,C \in \Obj(\mathcal{C})$ aus einer Abbildung: 
		\[
			\Mor_\mathcal{C}(B,C) \times \Mor_\mathcal{C}(A,B) \to \Mor_\mathcal{C}(A,C) \enspace, \quad (f,g) \mapsto f \circ g
		\]
		genannt die \Index{Komposition} in $\mathcal{C}$. 
	\end{enumerate}
	Dabei müssen folgende Axiome  erfüllt sein:
	\begin{enumerate}[(i)]
		\item Die Komposition ist \Index{assoziativ}: Für Objekte $A,B,C,D$ von $\mathcal{C}$ und $f \in \Mor_\mathcal{C}(C,D)$, $ g \in \Mor_\mathcal{C}(B,C)$ und $ h \in \Mor_\mathcal{C}(A,B)$ gilt immer
		\[
			f \circ (g \circ h) = (f \circ g) \circ h
		\]
		\item Die Komposition ist \Index{unital}: Für jedes Objekt $A$ von $\mathcal{C}$ gibt es einen Morphismus $\id_A \in \Mor_\mathcal{C}(A,A)$ so, dass
		\[
			\forall f \in \Mor_\mathcal{C}(B,A) : \id_A \circ f = f  \qquad  \text{und} \qquad  \forall f \in \Mor_\mathcal{C}(A,B) : f \circ \id_A = f
		\]
	\end{enumerate}
\end{definition}

\begin{beispiel}[{name=[Kategorien]}]
	\leavevmode
	\begin{enumerate}[(1),itemsep=0pt]
		\item Die Kategorie der Mengen $\SET$: Objekte sind Mengen und Morphismen sind Abbildungen.
		\item Die Kategorie der Gruppen $\GRUPPEN$: Objekte sind Gruppen und Morphismen sind Gruppenhomomorphismen.
		\item Die Kategorie der $K$-Vektorräume $K\text{-}\VR$: Objekte sind $K$-Vektorräume und die Morphismen sind $K$-lineare Abbildungen.
		\item Die Kategorie der $R$-Moduln $R\text{-}\MOD$: Objekte sind $R$-Moduln und Morphismen sind $R$-lineare Abbildungen.
		\item Die Kategorie der $C^\infty$-Mannigfaltigkeiten $C^\infty\text{-}\MAN$: Objekte sind $C^\infty$-Mannigfaltigkeiten und Morphismen sind $C^\infty$-Abbildungen.
		\item Die Kategorie der topologischen Räume $\TOP$: Objekte sind topologische Räume und Morphismen sind stetige Abbildungen.
		\item Die Kategorie der punktierten topologischen Räume $\TOP_\bullet$: Objekte sind punktiert"=topologische Räume und Morphismen sind punktiert-stetige Abbildungen.
		\item Die Kategorie $\HTOP$: Objekte sind topologische Räume und Morphismen sind Homotopieklassen von stetigen Abbildungen.
		\item Die Kategorie $\HTOP_\bullet$: Objekte sind punktiert-topologische Räume und Morphismen sind Homotopieklassen von punktiert-stetigen Abbildungen.
		\item Sei $G$ eine Gruppe. Wir erhalten eine Kategorie $\mathcal{C}_G$ mit genau einem Objekt $*$ und Morphismen $\Mor_{\mathcal{C}_G}(*,*) = G$. Die Komposition wird durch die
		Verknüpfung in der Gruppe festgelegt.
		\item Ist $\mathcal{C}$ eine Kategorie, so ist $\mathcal{C}^{\op}$ eine Kategorie, wobei $\Obj(\mathcal{C}^\op) = \Obj(\mathcal{C})$ und 
		\(
			\Mor_{\mathcal{C}^\op}(A,B) := \Mor_\mathcal{C}(B,A)
		\).
		Die Komposition ist gegeben durch $f \circ^\op g := g \circ f$.
	\end{enumerate}
\end{beispiel}

\begin{bemerkung}[{name=[kleine Kategorie]}]
	\leavevmode
	\begin{enumerate}[(i)]
		\item Eine Kategorie heißt \bet{klein}\index{Kategorie!klein}, wenn ihre Objekte eine Menge bilden.
		\item Statt $f \in \Mor_\mathcal{C}(A,B)$ schreiben wir oft $f \colon A \to B$ oder $A \xrightarrow{f} B$.
	\end{enumerate}
\end{bemerkung}

\begin{definition}[{name=[Funktor]}]
	Seien $\mathcal{C}$ und $\mathcal{D}$ Kategorien: Ein \Index{Funktor} $F \colon \mathcal{C} \to \mathcal{D}$ ordnet jedem Objekt $C$ von $\mathcal{C}$ ein Objekt $F(C)$ von 
	$\mathcal{D}$ und ordnet jedem Morphismus $f : C \to C'$ in $\mathcal{C}$ einen Morphismus $F(f) \colon F(C) \to F(C')$ in $\mathcal{D}$ zu. Dabei muss gelten:
	\[
		F(f \circ g) = F(f) \circ F(g) \quad \text{und} \quad F(\id_C) = \id_{F(C)}
	\]
\end{definition}

\begin{beispiel}[{name=[Funktoren]}]
	\leavevmode
	\begin{enumerate}[(1)]
		\item Es gibt offensichtliche \enquote{Vergiss}-Funktoren:
		\begin{gather}
			K\text{-}\VR \longrightarrow  \SET\qquad  R\text{-}\MOD \longrightarrow  \SET \qquad \TOP \longrightarrow  \SET\\
			\TOP_\bullet \longrightarrow  \SET  \qquad 
			\TOP_\bullet \longrightarrow   \TOP
		\end{gather}
		\item Sei $\mathcal{C}$ eine Kategorie und $C$ ein Objekt von $\mathcal{C}$. 
		Der durch $C$ \bet{dargestellte Funktor}\index{Funktor!dargestellter} $F_C \colon \mathcal{C} \to \SET$ ist definiert durch 
		\[
			\begin{array}{ll}
				F_C(A) = \Mor_\mathcal{C}(C,A) &\text{ für } A \in  \Obj(\mathcal{C}) \\
				F_C(f) \colon \Mor_\mathcal{C}(C,A) \to \Mor_\mathcal{C}(C,A'), g \mapsto f \circ  g &\text{ für } f \colon A \to A' \text{ in } \Mor(A,A')
			\end{array}
		\]
		\item Die Fundamentalgruppe definiert einen Funktor
		\[
			\pi_1 \colon \HTOP_\bullet \longrightarrow \GRUPPEN
		\]
	\end{enumerate}
	Ist $F \colon \mathcal{C} \to \mathcal{D}$ ein Funktor so schreiben wir oft kürzer und ungenauer $f_* := F(f)$.
	Kleine Kategorien und Funktoren bilden die Kategorie $\KAT$.
\end{beispiel}

\begin{definition}[{name=[natürliche Transformation]}]
	Seien $F,G \colon \mathcal{C} \to \mathcal{D}$ zwei Funktoren. 
	Eine \Index{natürliche Transformation} $\tau \colon F \to G$ ordnet jedem $C \in \Obj(\mathcal{C})$ einen Morphismus $\tau_C \colon F(C) \to G(C)$ in $\mathcal{D}$ zu, sodass für jedes $f \colon C \to C'$ in $\mathcal{C}$
	\[
		\begin{tikzcd}[sep=large]
			F(C) \rar["F(f)"] \dar["\tau_C"]& F(C') \dar["\tau_{C'}"] \\
			G(C) \rar["G(f)"] & G(C')
		\end{tikzcd}
	\]
	kommutiert.
\end{definition}

Natürliche Transformationen lassen sich komponieren. 
Für $\tau \colon F \to G$, $\eta \colon G \to H$ ist $\eta \circ \tau \colon F \to H$ gegeben durch
\marginnote{$F,G,H \colon \mathcal{C} \to \mathcal{D}$}
\[
	(\eta \circ \tau)_C :=  \eta_C \circ \tau_C \colon F(C) \to H(C)
\]
Genauer: Für feste kleine Kategorien $\mathcal{C}$ und $\mathcal{D}$ bilden die Funktoren $\mathcal{C} \to \mathcal{D}$ mit den natürlichen Transformationen eine Kategorie $\FUN(\mathcal{C}, \mathcal{D})$.

In der algebraischen Topologie werden topologische Fragen, wie zum Beispiel, wann $\mathbb{R}^n$ und $\mathbb{R}^m$ homöomorph sind, in algebraische Fragen übersetzt. 
Eine Möglichkeit für eine solche Übersetzung sind Funktoren von einer Kategorie von topologischen Räumen, z.B. $\TOP, \TOP_\bullet, \HTOP_\bullet$, in eine algebraisch Kategorie, z.B. $\GRUPPEN, K\text{-}\VR, \ABELGRUPPEN, R\text{-}\MOD$. 
Ein Beispiel für einen solchen Funktor ist die Fundamentalgruppe $\pi_1$. 
Ein Nachteil der Fundamentalgruppe ist, dass diese oft schwierig zu berechnen ist. 
Wir werden in dieser Vorlesung weitere Funktoren und Methoden für ihre Berechnung kennenlernen.

\newpage

\section{Kombinatorische Beschreibungen topologischer Räume} % (fold)
\label{sec:2}

\begin{beispiel}[{name=[kombinatorische Beschreibung eines Raumes]}]
	\Index{Graphen} sind einerseits topologische Räume und andererseits kombinatorische Objekte:
	\begin{align*}
		\vcenter{\hbox{\begin{tikzpicture}[scale=1.2]
			\coordinate (A) at (0,1);
			\coordinate (B) at (1,1);
			\coordinate (C) at (2,1);
			\coordinate (1) at (0,0);
			\coordinate (2) at (1,0);
			\coordinate (3) at (2,0);
			\foreach \x in {1,2,3} {
				\foreach \y in {A,B,C}{
					\draw[thick] (\x) -- (\y);
				};
			}
			\foreach \x in {A,B,C} {
				\foreach \y in {1,2,3}{
					\draw[thick] (\x) -- (\y);
				};
				\draw[fill=DodgerBlue3, thick] (\x) circle[radius=0.05] node[above]{$\x$};
			}
			\foreach \x in {1,2,3} {
				\draw[fill=DodgerBlue3,thick] (\x) circle[radius=0.05] node[below]{$\x$};
			}
		\end{tikzpicture} }}\qquad 
		&\longleftrightarrow \quad \begin{array}{l}
			V =  \set{A,B,C,1,2,3} \\
			E =  \set{ \rule{0cm}{0.3cm} { \scriptstyle \set{A,1}, \set{A,2}, \set{A,3} , \set{B,1}, \set{B,2}, \set{B,3}, \set{C,1} , \set{C,2}, \set{C,3}}}  
		\end{array}\\
		\vcenter{\hbox{\begin{tikzpicture}[scale=1.2]
			\coordinate (1) at (0,1);
			\coordinate (2) at (0.5,1.5);
			\coordinate (3) at (1,1);
			\coordinate (4) at (1,0);
			\coordinate (5) at (0,0);
			\coordinate (6) at (0.5,0.5);
			\foreach \x in {2,3,5,6} {
				\draw[thick] (1) -- (\x);
			}
			\foreach \x/\y in {2/3,3/4,3/6,4/5,4/6,5/6}{
				\draw[thick] (\x) -- (\y);
			}
			\foreach \x/\y in {1/left,2/above,3/right,4/right,5/left,6/right}{
				\draw[fill=DodgerBlue3, thick] (\x) circle[radius=0.05] node[\y]{$\x$};
			}
		\end{tikzpicture} }}\qquad
		&\longleftrightarrow \quad \begin{array}{l}
			V = \set{1,2,3,4,5,6} \\
			E = \set{ \rule{0cm}{0.3cm} {\scriptstyle \set{1,2}, \set{1,3}, \set{1,5} , \set{1,6}, \set{2,3}, \set{3,4}, \set{3,6} , \set{4,5}, \set{4,6}, \set{5,6}} }  
		\end{array}
	\end{align*}
\end{beispiel}

\begin{definition}[{name=[simplizialer Komplex]}]
	Ein \Index{simplizialer Komplex} $K=(V, \Sigma)$ besteht aus einer Menge $V$ und einer Menge $\Sigma$ von nichtleeren, endlichen Teilmengen von $V$, sodass gilt
	\begin{enumerate}[(i)]
		\item $\set{v} \in \Sigma$ für alle $v \in V$
		\item $\sigma \in \Sigma$, $\emptyset \not= \tau \subseteq \sigma \Longrightarrow \tau \in\Sigma.$ \hfill {\footnotesize(Abgeschlossen bzgl. Teilmengenbildung)}
	\end{enumerate}
	Die Elemente von $V$ heißen die \Index{Ecken} oder \Index{Vertices} von $K$. 
	Die Elemente von $\Sigma$ heißen die \Index{Simplizes} von $K$. 
	Enthält $\sigma \in \Sigma$ genau $n+1$ Elemente, so heißt $\sigma$ ein \bet{$\mathbf{n}$-Simplex}\index{n-Simplex@$n$-Simplex}. 
	Ist $\tau \subseteq \sigma$, wobei $\sigma$ ein $n$-Simplex und $\tau$ ein $n-1$-Simplex ist, so heißt $\tau$ eine \Index{Seite} von $\sigma$. 
	Ist $V$ geordnet, so heißt $K$ geordnet. Ist $V$ endlich, so heißt $K$ endlich.
\end{definition}

\begin{beispiel}
	\begin{align*}
		\vcenter{\hbox{\begin{tikzpicture}[scale=1.]
			\coordinate (0) at (0,0);
			\coordinate (1) at (1,0.8);
			\coordinate (2) at (1,-0.7);
			\coordinate (3) at (2,0.2);
			\filldraw[DodgerBlue2!40] (0) -- (1) -- (2) -- cycle ;
			\foreach \x/\y in {0/1,0/2,1/2,1/3,2/3} {
				\draw[thick] (\x) -- (\y);
			};
			\foreach \x/\pos in {0/left,1/above,2/below,3/right} {
				\draw[fill=DodgerBlue3,thick] (\x) circle[radius=0.06] node[\pos]{$\x$};
			}
		\end{tikzpicture}}} \qquad 
		\begin{array}{l}
			\minwidthbox{V}{1em} =  \set{0,1,2,3} \\
			\minwidthbox{\Sigma}{1em} =  \set{ \rule{0cm}{0.3cm} {\set{0}, \set{1}, \set{2}, \set{3}, \set{0,1}, \set{0,2},\set{1,2},\set{1,3},\set{2,3},\set{0,1,2}}}  
		\end{array}
	\end{align*}
	Ist $\sigma$ eine endliche Menge, so heißt $\Delta^\sigma := \enbrace[\big]{\sigma, \mathcal{P}(\sigma) \setminus \set{\emptyset}}$ der $\sigma$-Simplex. Für 
	$\sigma = \set{0, \ldots , n}$ schreiben wir $\Delta^n := \Delta^{\set{0, \ldots ,n}}$. Es ist
	\[
		\Delta^0  = \vcenter{\hbox{\begin{tikzpicture}
			\draw[fill=DodgerBlue3,thick] (0,0) circle[radius=0.06];
		\end{tikzpicture}}} , \qquad 
		\Delta^1 = \vcenter{\hbox{\begin{tikzpicture}
			\coordinate (0) at (0,0);
			\coordinate (1) at (0.4,0.3);
			\draw[thick] (1) -- (0);
			\foreach \x in {0,1} {
				\draw[fill=DodgerBlue3,thick] (\x) circle[radius=0.06];
			}
		\end{tikzpicture}}} , \qquad 
		\Delta^2 = \vcenter{\hbox{\begin{tikzpicture}
			\coordinate (0) at (0,0);
			\coordinate (1) at (0.6,0);
			\coordinate (2) at (60:0.6);
			\draw[thick, fill=DodgerBlue2!40] (1) -- (0) -- (2) -- cycle;
			\foreach \x in {0,1,2} {
				\draw[fill=DodgerBlue3,thick] (\x) circle[radius=0.06];
			}
		\end{tikzpicture}}}, \qquad 
		\Delta^3 = 
		\vcenter{\hbox{\begin{tikzpicture}
				\coordinate (0) at (0,0);
				\coordinate (1) at (1,0);
				\coordinate (2) at (0.5,0,0.866);
				\coordinate (3) at (0.5,0.866,0.433);
				\foreach \x in {0,1,2,3}{
					\foreach \y in {0,1,2,3}{
						\foreach \z/\col in {0/DodgerBlue2,1/DodgerBlue1,2/DodgerBlue3,3/DodgerBlue2}{
						\draw[thick, fill=\col!40, fill opacity=0.4] (\x) -- (\y) -- (\z) -- cycle;
						};
					};
				};
				\foreach \x in {0,1,2,3} {
					\draw[fill=DodgerBlue3,thick] (\x) circle[radius=0.06];
				}
			\end{tikzpicture}}}
	\]
\end{beispiel}

\begin{definition}[{name=[Unterkomplex]}]
	Sei $K=(V, \Sigma)$ ein simplizialer Komplex. 
	Ein \Index{Unterkomplex} von $K$ ist ein simplizialer Komplex $K_0 = (V_0, \Sigma_0)$ mit $V_0 \subseteq V$ und $\Sigma_0 \subseteq \Sigma$.
	
	Sei $K=(V,\Sigma)$ ein endlicher simplizialer Komplex. 
	Dann ist $K$ ein Unterkomplex von $\Delta^V$.
\end{definition}

\begin{definition}[{name=[simpliziale Abbildung]}]
	Eine \Index{simpliziale Abbildung} $f \colon K_1=(V_1, \Sigma_1) \to K_2=(V_2, \Sigma_2)$ zwischen simplizialen Komplexen ist eine Abbildung $f \colon V_1 \to V_2$, sodass $f(\sigma_1) \in \Sigma_2$ für alle $\sigma_1 \in \Sigma_1$. 
	Simpliziale Komplexe zusammen mit simplizialen Abbildungen bilden eine Kategorie $\SIMP$.
\end{definition}

\begin{definition}[{name=[geometrische Realisierung]}]
	Sei $\sigma$ eine endliche Menge. 
	Sei $\mathbb{R}^\sigma = \prod_{v \in \sigma} \mathbb{R}$ mit der Produkttopologie. 
	Wir setzen nun 
	\[
		\abs*{\Delta^\sigma} := \set*{x=(x_v)_{v \in \sigma} \in \mathbb{R}^\sigma \given \sum_{v \in \sigma}x_v = 1 , x_v \in [0,1] \,\, \forall v \in \sigma}  
	\]
	Ist $\tau \subseteq \sigma$, so erhalten wir eine Abbildung $\iota_\tau^\sigma \colon \abs*{\Delta^\tau} \to \abs*{\Delta^\sigma}$ indem wir $(x_v)_{v \in \tau} \in \abs*{\Delta^\tau}$ durch $x_v=0$ für $v \in \sigma \setminus \tau$ zu $(x_v)_{v \in \sigma}$ auffüllen. 
	Ist $\eta \subseteq \tau \subseteq \sigma$ so gilt offensichtlich $\iota^\sigma_\eta = \iota_\tau^\sigma \circ  \iota_\eta^\tau$.
	Sei nun $K=(V, \Sigma)$ ein simplizialer Komplex. 
	Die \Index{geometrische Realisierung} $\abs*{K}$ von $K$ ist definiert als 
	\[
		\abs*{K} := \coprod_{\sigma \in \Sigma} \set{\sigma} \times \abs*{\Delta^\sigma} \Big/ \!\sim
	\] 
	wobei $\sim$ die durch $(\tau, x) \sim \enbrace[\big]{\sigma, \iota_\tau^\sigma(x)}$ für $\tau \subseteq \sigma$ erzeugte Äquivalenzrelation ist.\marginnote{Anschaulich \enquote{klebt} $\sim$ die geometrischen Realisierungen aller Einzelteile passend zusammen} 
	Versehen mit der Quotiententopologie ist $\abs*{K}$ ein topologischer Raum.
\end{definition}

\begin{itemize}
	\item In Aufgabe 4 von Blatt 1 haben wir gezeigt, dass $\abs*{K}$ homöomorph zu einem Unterraum $X$ von $\mathbb{R}^n$ mit $n:= \# V$ ist via $\benbrace[\big]{\enbrace*{\set{\sigma}, (x_{v_i})_{v_i \in \sigma} } } \mapsto \sum_{v_i \in \sigma} x_{v_i} \cdot e_i$.
	\item In Aufgabe 2 von Blatt 2 haben wir gezeigt, dass $\abs*{\,\cdot\,} $ einen Funktor $\SIMP \to \TOP$ definiert.
	Dazu definiert man eine Abbildung zunächst auf den zugrunde liegenden Räumen durch
	\[
		\enbrace*{\sigma, \sum\nolimits_{v \in \sigma} x_v \cdot v} \longmapsto \enbrace*{f(\sigma), \sum\nolimits_{v \in \sigma} x_v \cdot f(v)} 
	\]
	und wendet diese unter Benutzung der universellen Eigenschaft der Quotiententopologie auf Repräsentanten an.\footnote{hier wird eine formale Summe anstatt eines Tupels benutzt; da $\sigma$ stets eine endliche Menge ist, macht dies aber keinen Unterschied.}
\end{itemize}

\begin{beispiel}[{name=[Rand des $n$-Simplizes und der Torus]}]
	\leavevmode
	\begin{itemize}
		\item Sei $\partial \Delta^n := \enbrace[\big]{\set{0, \ldots ,n}, \mathcal{P}\enbrace[\big]{\set{0,\ldots ,n}} \setminus \set{\emptyset, \set{0, \ldots , n}}} = $ 
		\enquote{$\Delta^n \setminus \set{0,\ldots ,n}$}. $\partial \Delta^n$ heißt der \Index{Rand des $n$-Simplizes}. Es gilt
		\begin{align*}
			\abs*{\partial \Delta^1} &= \enspace\vcenter{\hbox{\begin{tikzpicture}
			\draw[fill] (0,0) circle[radius=0.04];
			\draw[fill] (0.155,0) circle[radius=0.04];
		\end{tikzpicture}}} \enspace\cong S^0\\ 
			\abs*{\partial \Delta^2} &=  \vcenter{\hbox{\begin{tikzpicture}
			\coordinate (0) at (0,0);
			\coordinate (1) at (0.6,0);
			\coordinate (2) at (60:0.6);
			\draw[thick] (1) -- (0) -- (2) -- cycle;
		\end{tikzpicture}}} \cong S^1 \\
			\abs*{\partial \Delta^3} &= \vcenter{\hbox{
			\begin{tikzpicture}
				\coordinate (0) at (0,0);
				\coordinate (1) at (1,0);
				\coordinate (2) at (0.5,0,0.866);
				\coordinate (3) at (0.5,0.866,0.433);
				\foreach \x in {0,1,2,3}{
					\foreach \y in {0,1,2,3}{
						\foreach \z/\col in {0/DodgerBlue2,1/DodgerBlue1,2/DodgerBlue3,3/DodgerBlue2}{
						\draw[thick, fill=\col!40, fill opacity=0.4] (\x) -- (\y) -- (\z) -- cycle;
						};
					};
				};
			\end{tikzpicture}}} \cong S^2 
		\end{align*}
		Allgemein gilt $\abs*{\partial \Delta^n} \cong S^{n-1}$, siehe Aufgabe 1 von Blatt 3. 
		\item $T^2 = S^1 \times S^1 \cong 
		\vcenter{\hbox{\begin{tikzpicture}[scale=0.3]
			\draw (-3.5,0) .. controls (-3.5,2) and (-1.5,2.5) .. (0,2.5);
			\draw[xscale=-1] (-3.5,0) .. controls (-3.5,2) and (-1.5,2.5) .. (0,2.5);
			\draw[rotate=180] (-3.5,0) .. controls (-3.5,2) and (-1.5,2.5) .. (0,2.5);
			\draw[yscale=-1] (-3.5,0) .. controls (-3.5,2) and (-1.5,2.5) .. (0,2.5);

			\draw (-2,.2) .. controls (-1.5,-0.3) and (-1,-0.5) .. (0,-.5) .. controls (1,-0.5) and (1.5,-0.3) .. (2,0.2);

			\draw (-1.75,0) .. controls (-1.5,0.3) and (-1,0.5) .. (0,.5) .. controls (1,0.5) and (1.5,0.3) .. (1.75,0);
		\end{tikzpicture}}} \cong 
		\abs*{\vcenter{\hbox{\begin{tikzpicture}[>=to, scale=0.8]
			\filldraw[color=DodgerBlue2!20,] (0,0) rectangle (3,3);
			\foreach \x in {0,1,2}{
				\foreach \y in {0,1,2}{
					\draw (\x,\y) -- ++ (1,1);
				};
			};
			\draw (0,0) grid (3,3);
			\draw[very thick] (0,0) -- (0,3) -- (3,3) -- (3,0) -- cycle;
			\draw[->>, thick] (0,1) -- (0,1.5);
			\draw[->, thick] (1,3) -- (1.5,3);
			\draw[->, thick] (1,0) -- (1.5,0);
			\draw[->>, thick] (3,1) -- (3,1.5);
			\clip (-0.2,-0.2) rectangle (3.2,3.2);
		\end{tikzpicture}}}} $
	\end{itemize}
\end{beispiel}

\begin{definition}[{name=[Polyeder]}]
	Ein topologischer Raum $X$ heißt ein \Index{Polyeder}, falls er homöomorph zur geometrischen Realisierung eines simplizialen Komplexes ist. 
\end{definition}

\begin{definition}[{name=[Euler-Charakteristik]}]
	Sei $K=(V,\Sigma)$ ein endlicher simplizialer Komplex. Sei 
	\[
		a_n := \# \set[\big]{\sigma \in \Sigma \given \sigma \text{ ist ein $n$-Simplex von $K$}}
	\] 
	Dann heißt $\chi(K) := \sum_{n \in \mathbb{N}} (-1)^n \cdot a_n $ die \Index{Euler-Charakteristik} von $K$.
\end{definition}

Die Eulercharakteristik lässt sich einfach durch Abzählen bestimmen: 
\begin{align*}
	\chi \enbrace*{\vcenter{\hbox{\begin{tikzpicture}
		\coordinate (0) at (0,0);
		\coordinate (1) at (0.6,0);
		\coordinate (2) at (60:0.6);
		\draw[thick] (1) -- (0) -- (2) -- cycle;
	\end{tikzpicture}}} } &= 3 -3 = 0 \\
	\chi \enbrace[\bigg]{\vcenter{\hbox{\begin{tikzpicture}
		\coordinate (0) at (0,0);
		\coordinate (1) at (0.6,0);
		\coordinate (2) at (0.6,0.6);
		\coordinate (3) at (0,0.6);
		\draw[thick] (0) -- (1) -- (2) -- (3) -- cycle;
	\end{tikzpicture}}} } &= 4-4 =0 \\
	\chi \enbrace[\big]{\benbrace*{n\text{-Eck}}} &= n-n = 0
\end{align*}
Selbstverständlich funktioniert dies auch bei dreidimensionalen Komplexen:
\begin{align*}
	\chi \enbrace*{\partial \Delta^3 = \vcenter{\hbox{\begin{tikzpicture}
			\coordinate (0) at (0,0);
			\coordinate (1) at (1,0);
			\coordinate (2) at (0.5,0,0.866);
			\coordinate (3) at (0.5,0.866,0.433);
			\draw[sharp corners, dashed] (0) -- (1);
			\draw[sharp corners] (1) -- (2) -- (0) -- (3) -- (2) -- (1) -- (3) -- (0);
		\end{tikzpicture}}} } &= 4 - 6 +4 = 2 \\
	\chi\enbrace[\Bigg]{ \vcenter{\hbox{\begin{tikzpicture}
		\coordinate (0) at (0,0,-0.6);
		\coordinate (1) at (0.6,0,-0.6);
		\coordinate (2) at (0.6,0.6,-0.6);
		\coordinate (3) at (0,0.6,-0.6);
		\draw[thick, dashed] (3) -- (0) -- (1);
		\draw[thick] (1) -- (2) -- (3);
		\draw[thick, dashed] (0,0) -- (0);
		\draw[thick] (0.6,0) -- (0.6,0,-0.6);
		\draw[thick] (0.6,0.6) -- (0.6,0.6,-0.6);
		\draw[thick] (0,0.6) -- (0,0.6,-0.6) -- ++ (0.2,0);
		\draw[thick] (0,0) rectangle (0.6,0.6);
	\end{tikzpicture}}} } &= 8 -12 +6 = 2 \\
	\chi \enbrace*{ \vcenter{\hbox{\begin{tikzpicture}
		\coordinate (A1) at (0,0);
		\coordinate (A2) at (0.6,0.2);
		\coordinate (A3) at (1,0);
		\coordinate (A4) at (0.4,-0.2);
		\coordinate (B1) at (0.5,0.5);
		\coordinate (B2) at (0.5,-0.5);
		\begin{scope}[thick,dashed,]
		\draw (A1) -- (A2) -- (A3);
		\draw (B1) -- (A2) -- (B2);
		\end{scope}
		\begin{scope}[thick]
		\draw[] (A1) -- (A4) -- (B1);
		\draw[] (A1) -- (A4) -- (B2);
		\draw[] (A3) -- (A4) -- (B1);
		\draw[] (A3) -- (A4) -- (B2);
		\draw (B1) -- (A1) -- (B2) -- (A3) --cycle;
		\end{scope}
	\end{tikzpicture}}} } &= 6- 12  +8 = 2 
\end{align*}

\begin{satz}
	Seinen $K$ und $K'$ endliche simpliziale Komplexe. 
	Gilt $\abs*{K} \cong \abs*{K'}$, so gilt 
	\[
		\chi(K)= \chi(K').
	\]
\end{satz}
\begin{beweis}
	% \emph{Siehe \ref{sub:1210}.}
	\todo{auskommentieren}
\end{beweis}

Die Euler-Charakteristik ist also eine topologische Invariante von simplizialen Komplexen. 
Sie hängt nur von der topologischen Struktur von $\abs*{K}$ und nicht von der kombinatorischen Struktur von $K$ ab.

\newpage

\section{Simpliziale Homologie} % (fold)
\label{sec:3}

\begin{definition}[{name=[freier Modul]}]
	Sei $S$ eine Menge. Sei $R$ ein Ring.\marginnote{Ring mit Eins} 
	Der \bet{freie $R$-Modul mit Basis $S$}\index{freier Modul}, $R[S]$, besteht aus allen endlichen formalen $R$-Linearkombinationen
	\[
		\sum_{s \in S} r_s \cdot s \quad , r_s \in R \text{ und } r_s \not=0 \text{ für endliche viele }s \in S
	\]
	Die $R$-Modulstruktur auf $R[S]$ ist definiert durch:
	\begin{align*}
		r \cdot \enbrace*{\sum_{s\in S} r_s \cdot s} := \sum_{s \in S} (r \cdot r_s) \cdot s \qquad \qquad 
		\enbrace*{\sum_{s \in S} r_s \cdot s} + \enbrace*{\sum_{s \in S} r_s' s} := \sum_{s \in S} (r_s + r_s') \cdot s
	\end{align*}
	Mittels $s= \sum_{s' \in S} \delta_{s,s'} \cdot s'$ fassen wir $S$ als Teilmenge von $R[S]$ auf.
\end{definition}

Die freien Moduln sind durch die folgende universelle Eigenschaft charakterisiert:

\begin{lemma}[{name=[Universelle Eigenschaft von freien Moduln]},label=univ_frei_mod]
	Sei $S$ eine Menge, $R$ ein Ring. Es gilt\\
	\begin{minipage}[c]{0.7\textwidth}
		\begin{enumerate}[1)]
			\item Ist $M$ ein $R$-Modul und $i \colon S \to M$ eine Abbildung, so gibt es genau eine $R$-lineare Abbildung $\varphi \colon R[S] \to M$, die $i$ fortsetzt.
			\item Sei $f\colon R[S]\to M$ $R$-linear und $p\colon N\twoheadrightarrow M$ $R$-linear und surjektiv. Dann gibt es $\hat{f}\colon R[S]\to N$ mit $p\circ \hat{f}=f$.
		
			Moduln mit dieser Eigenschaft bezeichnet man als \bet{projektiv}\index{projektiver Modul}. Freie Moduln sind also projektiv.
		\end{enumerate}
	\end{minipage}\hfill
	\begin{minipage}[c]{0.25\textwidth}
		\[
			\begin{tikzcd}
				\mbox{ } & N \dar["p", two heads] \\
				R[S] \rar["f"] \urar["\hat{f}", dashed] & M
			\end{tikzcd}
		\]
	\end{minipage}
\end{lemma}
\begin{beweis}
	\leavevmode
	\begin{enumerate}[1)]
		\item Eine solche Abbildung ist gegeben durch
		\(
			\varphi \enbrace*{\sum_{s \in S} r_s \cdot s} = \sum_{s \in S} r_s \cdot i(s)
		\).
		Ist $\varphi'$ eine zweite, so gilt 
		\begin{align*}
			\hat{\varphi} \enbrace*{\sum_{s \in S} r_s \cdot s}  = \sum_{s \in S} r_s \cdot \hat{\varphi}(s) = \sum_{s  \in S} r_s \cdot i(s)
		\end{align*}
		\item Wähle für jedes $s \in S$ ein Urbild $j(s)\in N$ für $f(s) \in M$ unter $p\colon N \twoheadrightarrow M$. 
		Nun wende 1) auf $j \colon S \to N$ an, um $\hat{f} \colon R[S] \to N$ mit $\hat{f}(s) = j(s)$ zu erhalten. 
		Nun ist $ p \circ \hat{f}(s) = f(s)$ für alle $s \in S$. 
		Mit der Eindeutigkeit aus 1) folgt $p \circ \hat{f} = f$.\qedhere
	\end{enumerate}
\end{beweis}

\begin{beispiel}[{name=[freie Moduln]}]
	Viele uns bereits bekannte Strukturen sind freie Moduln:
	\begin{enumerate}[(i)]
		\item Sei $K$ ein Körper und $V,M,N$ seien $K$-Vektorräume, $f\colon V \to M$, $p\colon N \to M$ seien $K$-linear und $p$ surjektiv. 
		Dann gibt es $\hat{f}\colon V\to N$ mit $p \circ \hat{f} = f$. 
		Wir können Lemma \ref{univ_frei_mod} benutzen, da $V$ eine Basis $B$ hat, also $V \cong K[B]$.
		\item Sei $R=\mathbb{Z}$, sei $V= \nicefrac{\mathbb{Z}}{2 \mathbb{Z}}$ und $M = \nicefrac{\mathbb{Z}}{2 \mathbb{Z}}$, sowie $f= \id \colon V \to M$, $N= \mathbb{Z}$ und $p \colon \mathbb{Z} \twoheadrightarrow \nicefrac{\mathbb{Z}}{2 \mathbb{Z}}$ die Projektion mit $p(n)= n + 2 \mathbb{Z}$. 
		Dann gibt es \emph{keine} $\mathbb{Z}$-lineare Abbildung $\hat{f} \colon \nicefrac{\mathbb{Z}}{2 \mathbb{Z}} \to \mathbb{Z}$ mit $p \circ \hat{f} = f$, denn falls $\hat{f} (1+ 2 \mathbb{Z}) = n \in \mathbb{Z}$ ist, so folgt 
		\[
			2 n = 2 \cdot \hat{f}(1+ 2 \mathbb{Z}) = \hat{f} \enbrace[\big]{2 \cdot (1+ 2 \mathbb{Z})} = \hat{f} (0+ 2 \mathbb{Z}) =0
		\]
		also $n=0$ und somit $p \circ \hat{f}=0$. 
		$\nicefrac{\mathbb{Z}}{2 \mathbb{Z}}$ ist also kein freier Modul.
	\end{enumerate}
\end{beispiel}

\begin{definition}[{name=[{$n$-ter Kettenmodul}]}]
	Sei $K=(V, \Sigma)$ ein simplizialer Komplex. Sei $\Sigma_n := \set{\sigma \in \Sigma \given \sigma \text{ ist ein $n$-Simplex}}$. Sei $R$ ein Ring. Der 
	\bet{$\mathbf{n}$-te Kettenmodul}\index{n-ter Kettenmodul@$n$-ter Kettenmodul!simplizial} von $K$ über $R$ ist definiert als 
	\[
		C_n(K;R) := R[\Sigma_n]
	\]
	Ist $R=\mathbb{Z}$, so schreiben wir auch kurz $C_n(K) := C_n(K;\mathbb{Z})$.
\end{definition}

\begin{definition}[{name=[{Orientierung eines $(n-1)$-Simplizes als Teilmenge eines $n$-Simplizes}]}]
	Sei $K=(V,\Sigma)$ ein geordneter simplizialer Komplex. 
	Sei $\sigma = \set{v_0, \ldots ,v_n}$ ein $n$-Simplex von $K$. 
	Bezüglich der Ordnung von $K$ sei dabei $v_0 < v_1 < \ldots < v_n$. 
	Sei $\tau$ ein $(n-1)$-Simplex von $K$. Definiere
	\[
		\varepsilon_\sigma^\tau := \begin{cases}
			(-1)^i, &\text{ falls }\tau = \set{v_0, \ldots , v_{i-1}, v_{i+1}, \ldots , v_n} \\
			0 , &\text{ sonst}
		\end{cases}
	\]
\end{definition}

Es gilt offensichtlich $\varepsilon^\tau_\sigma \neq 0$ genau dann, wenn $\tau$ eine Seite von $\sigma$ ist.
Für den folgenden geordneten simplizialen Komplex gilt
\[
	\vcenter{\hbox{\begin{tikzpicture}
		\coordinate (0) at (0,0);
		\coordinate (1) at (0.5,0.866);
		\coordinate (2) at (1,0);
		\draw[very thick,
			fill=DodgerBlue2!40,
			postaction={decorate},
			decoration={markings,
				mark=at position 1/6 with {\arrow{to}},
				mark=at position 0.5 with {\arrow{to}},
				mark=at position 5/6 with {\arrow{to reversed}}}] 
		(0) -- (1) -- (2) -- cycle;
		\foreach \x/\pos in{0/left,1/above,2/right}{
			\node[\pos] at (\x) {$\x$};
		}
	\end{tikzpicture}}} \qquad 
	\varepsilon_{\set{0,1,2}}^{\set{0,1}} = 1 , \quad \varepsilon_{\set{0,1,2}}^{\set{0,2}}=-1 , \quad \varepsilon_{\set{0,1,2}}^{\set{1,2}} =1    
\]

\begin{definition}[{name=[{$n$-te Randabbildung}]}]
	Sei $K=(V,\Sigma)$ ein geordneter simplizialer Komplex. 
	Die \bet{$\mathbf{n}$-te Randabbildung}\index{n-te Randabbildung@$n$-te Randabbildung!simplizial} 
	\[
		\partial_n \colon \underset{=R[\Sigma_n]}{C_n(K;R)} \longrightarrow \underset{=R[\Sigma_{n-1}]}{C_{n-1}(K;R)}
	\]
	ist für ein Basiselement $\sigma\in \Sigma_n$ definiert durch
	\[
		\partial_n(\sigma) = \sum_{\mathclap{\tau \in \Sigma_{n-1}}} \varepsilon_\sigma^\tau \cdot \tau
	\]
\end{definition}

\begin{proposition}[{name=[{Für die Randabbildung gilt $\partial_{n-1} \circ \partial_n = 0$}]}]
	Es gilt
	\[
		\partial_{n-1} \circ \partial_n = 0
	\]
	Damit folgt $\im \partial_n \subseteq \ker \partial_{n-1}$.
\end{proposition}
\begin{beweis}
	Sei $\sigma$ ein $n$-Simplex von $K$. Dann gilt
	\begin{align*}
		\partial_{n-1} \enbrace[\big]{\partial_n (\sigma)} = \partial_{n-1} \enbrace*{ \smashoperator[r]{\sum_{\tau \in \Sigma_{n-1}}} \varepsilon_\sigma^\tau \cdot \tau} = 
		\sum_{\mathclap{\tau \in \Sigma_{n-1}}} \varepsilon_{\sigma}^\tau \cdot \partial_{n-1} (\tau) 
		 &= \smashoperator[r]{\sum_{\tau \in \Sigma_{n-1}}} \varepsilon_\sigma^\tau \enbrace*{ \sum_{\eta \in \Sigma_{n-2}} \varepsilon_\tau^\eta \cdot \eta} \\
		 &= \sum_{\eta \in \Sigma_{n-2}} \enbrace*{ \sum_{\tau \in \Sigma_{n-1}, \eta \subseteq \tau \subseteq \sigma} \!\!\!\!\! \! \! \varepsilon_\tau^\eta \cdot \varepsilon_\sigma^\tau}
		 \eta  \\
		 &\stackrel{\mathclap{\text{Lemma \ref{lem:38}}}}{=} \quad \enspace 0\qedhere
	\end{align*}
\end{beweis}

\begin{lemma}[label=lem:38,{name=[{Ein $(n-2)$-Simplex ist Seite von genau zwei $(n-1)$-Simplizes}]}]
	Sei $\sigma$ ein $n$-Simplex. Sei $\eta \subseteq \sigma$ ein $(n-2)$-Simplex. Dann gibt es genau zwei $(n-1)$-Simplizes $\tau, \tau'$ von $K$, die eine Seite von $\sigma$
	sind und $\eta$ als Seite enthalten. Es gilt
	\[
		\varepsilon_\tau^\eta \cdot \varepsilon^\tau_\sigma = - \varepsilon_{\tau'}^\eta \cdot \varepsilon^{\tau'}_\sigma
	\]
\end{lemma}
\begin{beweis}
	Sei $\sigma= \set{v_0, \ldots ,v_n}$ mit $v_0 < v_1 < \ldots < v_n$. 
	Dann ist $\eta = \set{v_0, \ldots , v_n} \setminus \set{v_i, v_j}$ mit $i<j$. 
	Dann sind $\tau = \set{v_0, \ldots , v_{i-1}, v_{i+1}, \ldots , v_n}$ und $\tau' = \set{v_0, \ldots , v_{j-1}, v_{j+1}, \ldots , v_n}$ die gesuchten $(n-1)$-Simplizes und es gilt
	\[
		\varepsilon_\sigma^\tau = (-1)^i \enspace, \enspace \varepsilon_\sigma^{\tau'} = (-1)^j \enspace, \enspace \varepsilon^\eta_\tau = (-1)^{j-1} \enspace, \enspace 
		\varepsilon_{\tau'}^\eta = (-1)^i \qedhere
	\]
\end{beweis}

Wir vereinbaren die folgende Notation: Für $\sigma \in \Sigma_n$, $ \sigma= \set{v_0, \ldots ,v_n} $ mit $v_0 < \ldots < v_n$ schreiben wir 
\[
	\delta^j \sigma := \set{v_0 , \ldots , v_{j-1}, v_{j+1}, \ldots , v_n}  \in \Sigma_{n-1}
\]
Dann ist $\partial_n(\sigma) = \sum_{j=0}^{n} (-1)^j \cdot \delta^j \sigma$.

\begin{definition}[{name=[{$n$-ter Homologiemodul}]}]
	Sei $K$ ein geordneter simplizialer Komplex. Der \bet{$\mathbf{n}$-te Homologiemodul}\index{n-ter Homologiemodul@$n$-ter Homologiemodul!simplizial} von 
	$K$ über $R$ ist definiert als\marginnote{In Aufgabe 3 von Blatt 3 haben wir gezeigt, dass die Homologiemoduln für verschiedene Ordnungen auf der gleichen Menge isomorph sind.}
	\[
		H_n(K;R) := \frac{\ker \partial_n \colon C_n(K;R) \to C_{n-1}(K;R)}{\im \partial_{n+1} \colon C_{n+1}(K;R) \to C_n(K;R)} 
	\]
	Für $n=0$ interpretieren wir $\partial_0$ als die Nullabbildung. Daher gilt 
	\[
		H_0(K;R) = \frac{C_0(K;R)}{\im \partial_1 \colon C_1(K;R) \to C_0(K;R)} 
	\]
\end{definition}

\begin{beispiel}
	Für ausreichend \enquote{einfache} Komplexe können wir die Homologie direkt ausrechnen:
	\begin{enumerate}[(1)]
		\item Sei $K= \bullet$. Dann ist $C_0(K;R) \cong R$ und $C_i(K;R) = 0$ für $i>0$. Weiter ist $H_0(K;R) = C_0(K;R) \cong R$ und $H_i(K;R) =0$ für $i>0$.
		\item Sein nun $K=K_n$ das $n$-Eck, wobei $n \ge 3$.  Also
		\begin{align*}
			\Sigma_0 &= \set[\big]{\tau_1=\set{1}, \tau_2=\set{2}, \ldots , \tau_n = \set{n} \rule{0cm}{0.9em}} \\
	 		\Sigma_1 &= \set[\big]{\sigma_1 = \set{1,2}, \ldots , \sigma_{n-1} = \set{n-1, n}, \sigma_n=\set{1,n} \rule{0cm}{0.9em}} 
		\end{align*}
		Dann gilt
		\begin{align*}
			C_0(K;R) &= R[\Sigma_0] \cong R^{n} \\
			C_1(K;R) &= R[\Sigma_1] \cong R^{n} \\
			C_i(K;R) &= 0 \text{ für $i>1$, insbesondere } H_i(K;R) = 0 \text{ für } i>1
		\end{align*}
		Es ist dann $\partial_2 = \partial_3 = \ldots  = 0$. Interessant ist $\partial_1 \colon C_1(K;R) \to C_0(K;R)$, denn es gilt
		\[
			\partial_1(\sigma_i) = \begin{cases}
				\tau_{i+1} - \tau_i, &\text{ falls }i=1, \ldots ,n-1\\
				\tau_n - \tau_1 , &\text{ falls } i=n
			\end{cases}
		\]
		Mit $\sigma_1' := \sigma_1, \sigma_2' := \sigma_2, \ldots , \sigma_{n-1}' := \sigma_{n-1}$ und $ \sigma_n' := - \sigma_n$ gilt dann
		\[
			\partial_1 (\sigma_i') = \tau_{i+1} - \tau_i
		\]
		mit der Konvention $\tau_{n+1} = \tau_1$. Also gilt
		\begin{align*}
			\partial_1 \enbrace*{\sum_{i=1}^{n} r_i \sigma_i' } = \sum_{i=1}^{n} r_i \enbrace*{\tau_{i+1}- \tau_i} = \sum_{i=1}^{n} (r_{i-1} - r_i) \tau_i  
		\end{align*}
		Es folgt
		\begin{align*}
			\ker \partial_1 &= \set*{\sum_{i=1}^{n} r_i \cdot \sigma_i' \given r_1=r_2= \ldots =r_n} \subseteq C_1(K;R) \\
			\im \partial_1 &= \set*{\sum_{i=1}^{n} s_i \cdot \tau_i \given \sum_{i=1}^{n} s_i =0} \subseteq C_0(K;R)
		\end{align*}
		Ist $\sum_{i=1}^{n} s_i =0$, so folgt für $r_n=0, r_1 = -s_1, r_2=-s_1 -s_2, \ldots , r_{n-1} = -s_1 - s_2 - \ldots - s_{n-1}$
		\[
			\partial_1 \enbrace*{\sum_{i=1}^{n} r_i \sigma_i'} = \sum_{i=1}^{n} s_i \tau_i 
		\]
		Es folgt $ %{\color{Firebrick1}[K] \in}\, 
		H_1(K;R) = \nicefrac{\ker \partial_1}{\im \partial_2} = \ker \partial_1 \cong R % \,{\color{Firebrick1}\ni 1}
		$ und 
		\[
			% { \color{Gold3} [\tau_i] =[p] \in }\,
			H_0(K;R) = \nicefrac{\ker \partial_0}{\im \partial_1} = \nicefrac{R[\Sigma_0]}{ \set*{\sum_{i=1}^{n} s_i \tau_i \given \sum_{i=1}^{n} s_i =0} } \cong R 
			% \, { \color{Gold3} \ni 1 }
		\]
		(via $\sum_{i=1}^{n} s_i \tau_i \longmapsto \sum_{i=1}^{n} s_i$ und Homomorphiesatz)
	\end{enumerate}
\end{beispiel}

\begin{definition}[{name=[{$R$-Kettenkomplex, $n$-te Homologie eines $R$-Kettenkomplexes}]}]
	Sei $R$ ein Ring. Ein \bet{$\mathbf{R}$-Kettenkomplex}\index{R-Kettenkomplex@$R$-Kettenkomplex} $(C_*, d_*)$ ist eine Folge von $R$-Moduln $(C_n)_{n \in \mathbb{N}}$ zusammen mit $R$-linearen Abbildungen $(d_n \colon C_n \to C_{n-1})_{n \ge 1}$, so dass $d_n \circ d_{n+1} = 0$. 
	Die $d_n$ heißen die \Index{Randabbildungen} von $(C_*,d_*)$, die $C_n$ die \Index{Kettenmoduln}.
	Die \bet{$\mathbf{n}$-te Homologie}\index{n-te Homologie@$n$-te Homologie} eines Kettenkomplexes $(C_*,d_*)$ ist definiert als
	\[
		H_n(C_*, d_*) := \frac{\ker d_n \colon C_n \to C_{n-1}}{\im d_{n+1} \colon C_{n+1} \to C_n} 
	\]
\end{definition}

\begin{bemerkung}
	\leavevmode
	\begin{enumerate}[(i)]
		\item Oft werden auch Kettenkomplexe betrachtet, die $\mathbb{Z}$ statt $\mathbb{N}$ verwenden.
		\item Ein Kettenkomplex $(C_*,d_*)$ heißt endlich erzeugt, wenn alle $C_n$ endlich erzeugte $R$-Moduln sind und $C_n \not= 0$ nur für endlich viele $n$ ist.
	\end{enumerate}
\end{bemerkung}

\begin{definition}[{name=[Simplizialer Kettenkomplex]}]
	Zu einem geordneten simplizialen Komplex $K$ heißt $\enbrace[\big]{C_*(K;R), \partial_*}$ der \bet{simpliziale Kettenkomplex}\index{simplizialer Kettenkomplex} von $K$ über $R$.
\end{definition}

Ist $K$ endlich, so ist $\enbrace[\big]{C_*(K;R), \partial_*}$ endlich erzeugt.


\newpage

\section{Euler-Charakteristik von Kettenkomplexen} % (fold)
\label{sec:4}

\begin{erinnerungA}[{name=[Torsionsgruppe]}]
	Sei $A$ eine abelsche Gruppe. Dann ist 
	\[
		T A := \set*{a \in A \given \exists n \ge 1 : n \cdot a = \underbracket[1pt][2pt]{a + \ldots + a}_{\text{$n$-mal}} = 0} 
	\]
	die \Index{Torsionsgruppe} von $A$. Ist $A$ endlich erzeugt, so ist auch $TA$ endlich erzeugt und es gibt Primzahlpotenzen $p_1^{n_1}, \ldots , p_k^{n_k}$ mit
	\marginnote{abelsche Gruppen sind $\mathbb{Z}$-Moduln}
	\[
		TA \cong \nicefrac{\mathbb{Z}}{(p_1^{n_1})} \oplus \ldots \oplus \nicefrac{\mathbb{Z}}{(p_k^{n_k})}
	\]
	Weiter gibt es dann $n$ mit $A \cong \mathbb{Z}^n \oplus TA$. $\Rg A := n$ ist der \bet{Rang} von $A$\index{Rang einer abelschen Gruppe}.
\end{erinnerungA}

\begin{lemma}[{name=[Rang einer endlich erzeugten abelschen Gruppe]}]
	Sei $A$ eine endlich erzeugte abelsche Gruppe. Dann gilt 
	\[
		\Rg A = \max \set[\big]{m \given \exists C \le A, C \cong \mathbb{Z}^m}.
	\]
\end{lemma}
\begin{beweis}
	Da $A \cong \mathbb{Z}^n \oplus T A$ ist, genügt zu zeigen: Ist $\varphi \colon \mathbb{Z}^m \to \mathbb{Z}^n \oplus T A$ ein injektiver Gruppenhomomorphismus, so ist 
	$m \le n$. Schreibe $\varphi = \varphi_0 \oplus \varphi_1$, also $\varphi(a) = \varphi_0(a) + \varphi_1(a)$ mit $\varphi_0 \colon \mathbb{Z}^m \to \mathbb{Z}^n$, 
	$\varphi_1 \colon \mathbb{Z}^m \to T A$. 

	Behauptung: $\varphi_0$ ist injektiv. Zu $v \in \mathbb{Z}^m$ wähle $k\ge 1$ mit $k \cdot \varphi_1(v) = 0 \in T A$. Dann ist $\varphi_1(k v) = k \varphi_1(v)=0$.
	Ist $v \not= 0$, so ist $k \cdot v \not= 0 \in \mathbb{Z}^m$, also ist $\varphi(k v) \not= 0$ und damit $\varphi_0(k v) \not= 0$. Es folgt $\varphi_0(v) \not= 0$, da sonst
	$\varphi_0(k \cdot v) = k \cdot \varphi_0(v)= 0$. Damit ist die Behauptung gezeigt.

	Gruppenhomomorphismen $\varphi_0 \colon \mathbb{Z}^m \to \mathbb{Z}^n$ werden durch $n \times m$-Matrizen beschrieben: Es gibt $A= (a_{ij}) \in \mathbb{Z}^{n \times m}$ mit
	\[
		\varphi_0 \begin{pmatrix}
			z_1 \\ \vdots \\ z_m
		\end{pmatrix} = \begin{pmatrix}
			\sum_{j=1}^{m} a_{1j} \cdot z_j \\
			\vdots \\
			\sum_{j=1}^{m} a_{nj} \cdot z_j
		\end{pmatrix}
	\]
	Ist $m >n$, so hat $A$, aufgefasst als Matrix über $\mathbb{Q}$, einen Kern. Es gibt also einen Vektor $w = \begin{psmallmatrix}a_1/b_1 \\ : \\ a_m/b_m\end{psmallmatrix}$
	mit $A w = 0$, $w \not=0$. Dann ist $(b_1, \ldots , b_m) \cdot w \in \mathbb{Z}^m$ und 
	\[
		\varphi_0 \enbrace[\big]{(b_1, \ldots , b_m)\cdot w} = A \cdot \enbrace[\big]{(b_1, \ldots , b_m) \cdot w} = (b_1, \ldots , b_m) \cdot A \cdot w = 0
	\]
	Dies ist ein Widerspruch zur Injektivität von $\varphi_0$. Also gilt $m \le n$.
\end{beweis}

\begin{definition}[{name=[Kurze exakte Sequenz]}]
	Seien $M_0, M_1, M_2$ $R$-Moduln und $f_0 \colon M_0 \to M_1$, $f_1 \colon M_1 \to M_2$ $R$-lineare Abbildungen. Dann heißt 
	\begin{equation*}
		\begin{tikzcd}
			M_0 \rar["f_0", hook] & M_1 \rar["f_1", two heads]& M_2 
		\end{tikzcd} \label{eq:43} \tag{*}
	\end{equation*}
	eine \Index{kurze exakte Sequenz}, wenn gilt:
	\begin{multicols}{3}
		\begin{enumerate}[(i)]
			\item $f_0$ ist injektiv,
			\item $\im f_0 = \ker f_1$,
			\item $f_1$ ist surjektiv.
		\end{enumerate}
	\end{multicols}
\end{definition}

Oft sagt man $M_0 \xrightarrow{f_0} M_1 \xrightarrow{f_1} M_2$ ist exakt in $M_1$, wenn $\im f_0 = \ker f_1$ ist. Dann ist \eqref{eq:43} eine kurze exakte Folge, wenn
\[
	\begin{tikzcd}
		0 \rar & M_0 \rar["f_0"] & M_1 \rar["f_1"]& M_2 \rar & 0
	\end{tikzcd}
\]
exakt in $M_0, M_1$ und $M_2$ ist.

\begin{beispiel}
	Es gibt zahlreiche Beispiele für kurze exakte Sequenzen:
	\begin{enumerate}[1)]
		\item \[
			\begin{tikzcd}[row sep=0.3em]
				M_0 \rar["i",hook] & M_0 \oplus M_1 \rar["p",two heads] & M_1 \\
				v_0 \rar["i", mapsto] & (v_0,0) \enspace (v_0,v_1) \rar["p",mapsto] & v_1
			\end{tikzcd}
		\]
		ist eine kurze exakte Folge.
		\item $\begin{tikzcd}
			\mathbb{Z} \rar[hook,"\cdot n"] & \mathbb{Z} \rar[two heads] &  \nicefrac{\mathbb{Z}}{n \mathbb{Z}}
		\end{tikzcd}$
		ist eine kurze exakte Folge.
	\end{enumerate}
\end{beispiel}

\begin{bemerkung}[{name=[Spaltung einer kurzen exakten Sequenz]}]
	Eine \Index{Spaltung} für eine kurze exakte Folge $M_0 \xrightarrow{f_0} M_1 \xrightarrow{f_1} M_2$ ist eine $R$-lineare Abbildung $s \colon M_2 \to M_1$ mit $f_1 \circ s = \id_{M_2}$. In diesem Fall erhalten wir einen Isomorphismus $M_0 \oplus M_2 \to M_1$, $(v_0, v_2) \mapsto f_0(v_0) + s(v_2)$.
	\begin{description}
		\item[Injektivität:] Sei $(v_0,v_2) \in M_0 \oplus M_2$ mit $f_0(v_0) + s(v_2) =0$. Dann gilt
		\[
			0 = f_1 \enbrace[\big]{f_0(v_0) + s(v_2)} = v_2 \implies f_0(v_0)=0 \implies v_0 =0
		\]
		\item[Surjektivität:] Sei $v_1 \in M_1$. Betrachte $v_2 := f_1(v_1)$. Dann ist $v_1 -s(v_2) \in \ker f_1 = \im f_0$, also gibt es $v_0$ in $M_0$ mit
		$f_0(v_0) = v_1 - s(v_1)$. Damit ist $v_1= f_0(v_0)+ s(v_1)$.
	\end{description}
\end{bemerkung}

\begin{bemerkung}[{name=[{Existenz einer Spaltung, wenn $M_2$ frei ist}]}]
	\begin{enumerate}[1)]
		\item Die kurze exakte Sequenz $\begin{tikzcd}
			\mathbb{Z} \rar[hook,"\cdot n"] & \mathbb{Z} \rar[two heads] &  \nicefrac{\mathbb{Z}}{n \mathbb{Z}}
		\end{tikzcd}$
		spaltet nicht.
		\item Ist $\begin{tikzcd}
			M_0 \rar["f_0"] & M_1 \rar["f_1"] & M_2
		\end{tikzcd}$
		eine kurze exakte Folge mit $M_2$ frei, also $M_2 \cong R[S]$, so spaltet die Folge.
		Insbesondere ist $M_1 \cong M_0 \oplus M_2$.
		Dazu wenden wir Lemma \ref{univ_frei_mod} an und erhalten eine $R$-lineare Abbildung $s \colon M_2 \to M_1$, sodass folgendes Diagramm kommutiert
		\[
			\begin{tikzcd}
				\mbox{ } & M_1 \dar["f_1",two heads] \\
				R[S] \urar["s", dashed] \rar["\id"] &  R[S]
			\end{tikzcd}
		\]
		$s$ ist also die gesuchte Spaltung.
		Dieses Vorgehen funktioniert bei jedem projektiven Modul.\index{projektiver Modul}
	\end{enumerate}
\end{bemerkung}

\begin{bemerkung}[{name=[Der Rang ist additiv]},label=bem:47]
	Der Rang für alle endlich erzeugten abelschen Gruppen ist additiv:
	\[
		\Rg(A \oplus B) =\Rg A + \Rg B,
	\]
	da $A \cong \mathbb{Z}^{\Rg A} \oplus T A$, $B \cong \mathbb{Z}^{\Rg B} \oplus TB$, also 
	\[
		A \oplus B \cong \mathbb{Z}^{\Rg A + \Rg B} \oplus \underbrace{T A \oplus T B}_{=T(A \oplus B)}
	\]
\end{bemerkung}

\begin{proposition}[{name=[{Der Rang ist für kurze exakte Folgen endl. erz., abelscher Gruppen additiv}]},label=prop:48]
	Der Rang von endlich erzeugten abelschen Gruppen ist additiv für kurze exakte Folgen: Ist
	\begin{equation*}
		\begin{tikzcd}
			A \rar["i",hook] & B \rar["p", two heads] & C
		\end{tikzcd} \label{eq:48}\tag{$\star$}
	\end{equation*}
	eine kurze exakte Folge von endlich 
	erzeugten abelschen Gruppen, so gilt $\Rg B = \Rg A + \Rg C$.
\end{proposition}
\begin{beweis}
	Spaltet \eqref{eq:48}, so ist $B \cong A \oplus C$ und die Behauptung folgt aus der \autoref{bem:47}.
	Wir können annehmen, dass $C = \mathbb{Z}^n \oplus T C$ gilt. Wir erhalten folgendes Diagramm: 
	\[
		\begin{tikzcd}
			A \rar["i",hook] \dar[equal] & p ^{-1}(\mathbb{Z}^n) \rar["p", two heads] \dar[hook,"i''"]& \mathbb{Z}^n \dar[hook,"i'"]\\
			A \rar["i",hook] & B \rar["p", two heads] \dar[two heads,"p''"]& C \dar[two heads,"p'"] & \text{\eqref{eq:48}}\\
			& \nicefrac{B}{p ^{-1}(\mathbb{Z}^n)} \rar[two heads, hook,"p' \circ p"]& T C
		\end{tikzcd} 
	\]
	Ist $p' \circ p(b)=0$, so gibt es $v \in \mathbb{Z}^n$ mit $i'(v)= p(b)$. Nun gibt es $v' \in p ^{-1}(\mathbb{Z}^n)$ mit $p(v')=v$. Es folgt 
	\[
		p \enbrace[\big]{i''(v')} = p(b) 
	\]
	also $i''(v') -b \in \ker p = i(A)$. Da $i(A) \subseteq i'' \enbrace*{p ^{-1} (\mathbb{Z}^n)}$, folgt $b \in i'' \enbrace*{p ^{-1}(\mathbb{Z}^n)}$.
	In dem Diagramm sind die Spalten und Zeilen exakt. Da $\mathbb{Z}^n$ frei ist, spaltet 
	$\begin{tikzcd}
		A \rar[hook,"i"] & p^{-1}(\mathbb{Z}^n) \rar["p", two heads] & \mathbb{Z}^n
	\end{tikzcd}$
	und es gilt 
	\[
		\Rg A + \Rg C = \Rg A + n = \Rg \enbrace*{p ^{-1}(\mathbb{Z}^n)} 
	\]
	Es bleibt zu zeigen: $\Rg B = \Rg  \enbrace*{p ^{-1}(\mathbb{Z}^n)}$. \medskip \\
	Nebenbei: Ist $A' \subseteq B'$, so gilt $\Rg A' \le \Rg B'$, denn 
	\[
		\Rg A' = \max \set*{n' \given \mathbb{Z}^{n'} \cong C' \le A'} \le \max \set*{n' \given \mathbb{Z}^{n'} \cong C' \le B'} = \Rg B'.
	\]
	Betrachte die kurze exakte Folge
	\[
		\begin{tikzcd}
			p^{-1}(\mathbb{Z}^n) \rar & B \rar & \nicefrac{B}{p ^{-1}(\mathbb{Z}^n)}
		\end{tikzcd}
	\]
	Wegen $\nicefrac{B}{p ^{-1}(\mathbb{Z}^n)} \cong T C$ gibt es ein
	$k$ mit $k \cdot \enbrace*{b+ p ^{-1}(\mathbb{Z}^n)} =0$ für alle $b \in B$. Also $k \cdot b \in p ^{-1}(\mathbb{Z}^n)$. Es folgt 
	$kB \subseteq  p ^{-1}(\mathbb{Z}^n)$. Ist  $B \cong \mathbb{Z}^{\Rg B}+ TB$, so ist $k B \cong (k \mathbb{Z})^{\Rg B}+ T(kB)$ und $\Rg kB = n = \Rg B$. 
	Mit der Nebenbemerkung folgt $\Rg p ^{-1}(\mathbb{Z}^n) = \Rg B$ aus $kB \subseteq  p ^{-1}(\mathbb{Z}^n) \subseteq B$.
\end{beweis}

Man kann die Proposition auch mittels Tensorprodukten beweisen.
Dazu geht man wie folgt vor:
\begin{itemize}
	\item $\mathbb{Q} \otimes_\mathbb{Z} \colon \mathbb{Z}\text{-}\MOD \to \mathbb{Q}\text{-}\VR $ ist ein Funktor.
	\item $\Rg A= \dim_\mathbb{Q} (\mathbb{Q} \otimes_\mathbb{Z} A)$
	\begin{enumerate}[(a)]
		\item $\mathbb{Q} \otimes_\mathbb{Z} (X \oplus Y) = \mathbb{Q} \otimes_\mathbb{Z} X \oplus  \mathbb{Q} \otimes_\mathbb{Z} Y$
		\item $\mathbb{Q} \otimes_\mathbb{Z} \mathbb{Z}^n \cong \mathbb{Q}^n$
		\item $\mathbb{Q} \otimes_\mathbb{Z} T \cong 0$ für $T$ ein Torsionsmodul.
	\end{enumerate}
	\item Ist $\begin{tikzcd}[cramped,sep=small]
		A \rar & B \rar & C
	\end{tikzcd}$ eine kurze exakte Folge, so ist auch 
	\(
		\begin{tikzcd}[cramped,sep=small]
				\mathbb{Q} \otimes_\mathbb{Z} A \rar & \mathbb{Q} \otimes_\mathbb{Z} B \rar & \mathbb{Q} \otimes_\mathbb{Z} C
			\end{tikzcd}
	\) 
	eine kurze exakte Folge.
	\item Jede kurze exakte Folge $\begin{tikzcd}
		V_0 \rar & V_1 \rar & V_2
	\end{tikzcd}$ von $\mathbb{Q}$-Vektorräumen spaltet und daher gilt $V_1 \cong V_0 \oplus V_2$ und $\dim_\mathbb{Q}(V \oplus W) = \dim_\mathbb{Q}(V)+ \dim_\mathbb{Q} (W)$.
\end{itemize}

\begin{definition}[{name=[{Eulercharakteristik von Z-Kettenkomplexen}]}]
	Sei $(C_*,d_*)$ ein endlich erzeugter $\mathbb{Z}$-Kettenkomplex.
	\[
		\chi(C_*) := \sum_{i=0}^{\infty} \Rg(C_i) = \sum_{i=0}^{n_0} \Rg(C_i) 
	\]
	heißt die \bet{Eulercharakteristik}\index{Euler-Charakteristik!von $\mathbb{Z}$-Kettenkomplexen} von $(C_*,d_*)$. 
	Dabei ist $n_0$ so gewählt, dass $C_n =0$ für alle $n >n_0$.
\end{definition}

\begin{satz}[{name=[Eulercharakteristik eines endlich erzeugten Z-Kettenkomplexes]},label=satz:410]
	Sei $(C_*,d_*)$ ein endlich erzeugter $\mathbb{Z}$-Kettenkomplex. Dann gilt 
	\[
		\chi(C_*) = \sum_{i=1}^{\infty} (-1)^i \Rg H_i(C_*,d_*).
	\]
\end{satz}
\begin{beweis}
	Sei $B_{n+1} := \im d_{n+1} \colon C_{n+1} \to C_n$ und $Z_n := \ker d_n \colon C_n \to C_{n-1}$. Also $H_n(C_*,d_*) = \nicefrac{Z_n}{B_{n+1}}$.
	Als Untermoduln von $C_n$ sind $Z_n$ und $B_n$ endlich erzeugt (LA2). Insbesondere ist auch $H_n(C_*,d_*)$ endlich erzeugt und der Rang somit definiert. Auch $B_{n+1}$
	ist Untermodul des endlich erzeugten $\mathbb{Z}$-Moduls $C_n$ und somit endlich erzeugt.
	Wir erhalten kurze exakte Folgen:
	\[
		\begin{tikzcd}[row sep=0.4em]
			B_{n+1} \rar[hook] & Z_n \rar[two heads] & H_n(C_*, d_*) \\
			Z_n \rar[hook] & C_n \rar[two heads] & B_n \quad 
		\end{tikzcd}
	\]
	Nach \autoref{prop:48} gilt $\Rg Z_n = \Rg B_{n+1} + \Rg H_n(C_*,d_*)$ und $\Rg C_n = \Rg Z_n + \Rg B_n$. Also gilt
	\begin{align*}
		\chi(C_*) = \sum_{i=0}^{\infty} (-1)^i \Rg C_i &= \sum_{i=0}^{\infty} (-1)^i \enbrace[\big]{\Rg Z_i + \Rg B_i} \\  
		&= \sum_{i=0}^{\infty} (-1)^i \enbrace[\Big]{\Rg H_i(C_*,d_*) + \Rg B_{i+1} + \Rg B_i} \\
		&= \sum_{i=0}^{\infty} (-1)^i \Rg H_i(C_*,d_*) + \sum_{i=0}^{\infty} (-1)^i \Rg B_{i+1} + \sum_{i=0}^{\infty}(-1)^i \Rg B_i \\
		&= \sum_{i=0}^{\infty} (-1)^i \Rg H_i(C_*,d_*)  \qedhere  
	\end{align*}
\end{beweis}

\begin{korollar}[{name=[Die Eulercharakteristik eines simplizialen Komplexes]}]
	Sei $K$ ein endlicher, geordneter simplizialer Kettenkomplex. 
	Dann gilt\marginnote{Eine vereinfachte Variante dieser Aussage kam in Aufgabe 2 von Blatt3 vor: Dort haben wirden Homologiemodul über einen Körper betrachtet. Deshalb konnte wir da dann die Dimensionsformeln anwenden, um die beiden Vorraussetzungen für die Rechnung in \ref{satz:410} 
	zu erhalten.}
	\[
		\chi(K) = \sum_{i=1}^{\dim K} (-1)^i \Rg H_i(K;\mathbb{Z})
	\]
\end{korollar}
\begin{beweis}
	Wende den \autoref{satz:410} auf den simplizialen Kettenkomplex von $K$ an. Da 
	\[
		\Rg C_n(K;\mathbb{Z}) = \Rg \mathbb{Z}[\Sigma_n] = \# \text{ $n$-Simplizes in $K$}
	\]
	ist
	\[
		\chi(K) = \chi \enbrace[\big]{C_*(K;\mathbb{Z}), \partial_*} = \sum_{i=0}^{\infty} (-1)^i \Rg H_i \enbrace[\big]{C_*(K;\mathbb{Z}), \partial_*}  =
		\sum_{i=0}^{\infty} (-1)^i \Rg H_i \enbrace*{K;\mathbb{Z}} \qedhere
	\]
\end{beweis}

\cleardoubleoddemptypage
\pagenumbering{Alph}
\setcounter{page}{1}
\printindex
\listoffigures
% \todototoc
% \listoftodos[To-do's und andere Baustellen]
\end{document}
